\documentclass[11pt,a4paper]{article}
\usepackage[margin=2.5cm]{geometry}
\usepackage{amsmath,amssymb}
\usepackage{bm}
\usepackage{booktabs}
\usepackage{float}
\usepackage{fancyhdr}
\usepackage[colorlinks=true,linkcolor=blue,citecolor=blue,urlcolor=blue]{hyperref}

\pagestyle{fancy}
\fancyhf{}
\rhead{HMC implementation, $N_f=4$}
\cfoot{\thepage}
\setlength{\headheight}{14pt}

%Ruled float for pseudocode (does not depend on the algorithm/algorithmicx packages)
\floatstyle{ruled}
\newfloat{algorithm}{tbp}{loa}
\floatname{algorithm}{Algorithm}

\newcommand{\n}{\mathbf{n}}
\newcommand{\z}{\mathbf{z}}
\newcommand{\Nf}{N_f}
\newcommand{\Npf}{N_{\mathrm{pf}}}
\newcommand{\NMD}{N_{\mathrm{MD}}}
\newcommand{\code}[1]{\texttt{#1}}
\newcommand{\hmcdoc}{\texttt{HMC\_doc.pdf}}

\title{\textbf{DOCUMENTATION}\\[0.8em]
\normalsize\emph{Hybrid Monte Carlo implementation for the four-flavor Schwinger model}}
\author{JFN Castellanos}
\date{}

\begin{document}
\maketitle

\noindent
This document describes the extension of the Hybrid Monte Carlo (HMC) code to four degenerate flavors of
Wilson fermions. It is a supplement to the documentation of the two-flavor code, \hmcdoc, and equation numbers
written as (H.$n$) refer to equation ($n$) of that document. All conventions are unchanged: the lattice and the
link variables, the $\gamma$ matrices of eq.~(H.15), the Wilson operator of eq.~(H.14), the antiperiodic boundary
conditions of the fermion fields in the time direction, and the definition of the derivative with respect to a link,
eq.~(H.23). Section~\ref{sec:action} derives the effective action, Sections~\ref{sec:heatbath}--\ref{sec:force}
give the pseudofermion heatbath, the Hamiltonian and the force, Section~\ref{sec:cost} discusses the cost and
alternative formulations, Section~\ref{sec:code} maps the equations onto the code, Section~\ref{sec:measurements}
comments on the measurements, and Section~\ref{sec:validation} reports the tests of the implementation.

The four-flavor version is selected at compile time,
\begin{verbatim}
    cmake -S . -B build -DFOUR_FLAVORS=ON
    cmake --build build
\end{verbatim}
and is only available for Wilson fermions: CMake stops with an error if \code{-DFOUR\_FLAVORS=ON} is combined with
\code{-DCLOVER=ON} or \code{-DTWISTED\_MASS=ON}. The runtime input parameters are the same as for the two-flavor
Wilson build.

%----------------------------------------------------------------------------------------------------------
\section{Effective action for four degenerate flavors}
\label{sec:action}

For $\Nf$ degenerate flavors without any flavor-mixing term, the Dirac operator in flavor space is
$\mathbb{1}_{\Nf}\otimes D$, with $D$ the Wilson operator of eq.~(H.14), and the partition function is
\begin{equation}
    Z = \int \mathcal{D}[U]\, e^{-S_G[U]}\, \bigl[\det(D)\bigr]^{\Nf},
    \qquad
    S_G[U] = \beta \sum_{\n} \mathrm{Re}\bigl(1-U_{01}(\n)\bigr).
    \label{eq:Z}
\end{equation}
The Wilson operator is $\gamma_5$-hermitian, $D^\dagger = \gamma_5 D \gamma_5$, eq.~(H.19). This implies
$\det(D) = \det(D^\dagger) = \det(D)^*$, i.e.\ the determinant is real (but not necessarily positive). For
$\Nf = 4$ we can therefore write
\begin{equation}
    \bigl[\det(D)\bigr]^4 = \bigl[\det(D)\det(D^\dagger)\bigr]^2 = \bigl[\det(DD^\dagger)\bigr]^2 \geq 0,
    \label{eq:det4}
\end{equation}
so the four-flavor weight is the square of the two-flavor weight and has no sign problem. Each factor
$\det(DD^\dagger)$ is represented by a Gaussian integral over a complex pseudofermion field, as in eq.~(H.18).
With two independent fields $\Phi_1$ and $\Phi_2$,
\begin{equation}
    \bigl[\det(DD^\dagger)\bigr]^2
    = \pi^{-2N} \int \mathcal{D}[\Phi_{1,R}]\,\mathcal{D}[\Phi_{1,I}]\,\mathcal{D}[\Phi_{2,R}]\,\mathcal{D}[\Phi_{2,I}]\,
      \exp\Bigl(-\sum_{k=1}^{2} \Phi_k^\dagger \bigl(DD^\dagger\bigr)^{-1} \Phi_k\Bigr),
    \label{eq:pf}
\end{equation}
where $N = 2N_xN_t$ is the number of complex components of each field. The effective action that replaces
eq.~(H.20) is therefore
\begin{equation}
    S[U,\Phi_1,\Phi_2] = \beta \sum_{\n} \mathrm{Re}\bigl(1-U_{01}(\n)\bigr)
    + \sum_{k=1}^{2} \sum_{\n,\n'} \sum_{\alpha,\beta=0}^{1}
      \Phi_{k,\alpha}^*(\n') \bigl[(DD^\dagger)^{-1}\bigr]^{\alpha\beta}_{\n',\n} \Phi_{k,\beta}(\n).
    \label{eq:S}
\end{equation}
Each pseudofermion field accounts for one pair of degenerate flavors. More generally, $\Nf = 2\Npf$ flavors are
simulated with $\Npf$ pseudofermion fields, and this is how the code is written: the number of fields is the
compile-time constant \code{N\_pf = Nf/2}, with \code{Nf = 4} if the code is compiled with \code{FOUR\_FLAVORS}
and \code{Nf = 2} otherwise. For $\Npf = 1$ all the expressions below reduce exactly to those of Section~2 of
\hmcdoc.

%----------------------------------------------------------------------------------------------------------
\section{Pseudofermion heatbath}
\label{sec:heatbath}

At the beginning of every trajectory both pseudofermion fields are generated anew. For $k = 1,2$ we draw a
complex field $\chi_k$ with probability $P[\chi_k] \propto \exp(-\chi_k^\dagger\chi_k)$, i.e.\ the real and
imaginary parts of every component are independent normal variables with zero mean and variance $1/2$, and set
\begin{equation}
    \Phi_k = D\,\chi_k, \qquad k = 1,2.
    \label{eq:heatbath}
\end{equation}
Since $\chi_k^\dagger\chi_k = \Phi_k^\dagger (D^\dagger)^{-1} D^{-1}\Phi_k = \Phi_k^\dagger (DD^\dagger)^{-1}\Phi_k$
and the Jacobian of the change of variables depends on $U$ only, at fixed $U$ the field $\Phi_k$ is distributed
according to $\exp\bigl(-\Phi_k^\dagger(DD^\dagger)^{-1}\Phi_k\bigr)$, as required by eq.~\eqref{eq:pf}. On the
configuration where they are generated, the fields satisfy
\begin{equation}
    \Phi_k^\dagger (DD^\dagger)^{-1}\Phi_k = \chi_k^\dagger\chi_k ,
    \label{eq:identity}
\end{equation}
which is a convenient check of the operators and of the solver.

The two fields must be drawn independently. If the same $\chi$ were used for both, then $\Phi_1 = \Phi_2 = \Phi$
and the fermion action would be $2\Phi^\dagger(DD^\dagger)^{-1}\Phi$. The Gaussian integral of this action is
proportional to $\det(DD^\dagger)$, i.e.\ to the \emph{two}-flavor weight. The simulation would then sample the
wrong theory while the force would still be the exact derivative of the action, so the error would not show up in
the energy violation or in the acceptance rate. In the code, a new $\chi_k$ is drawn for each $k$.

%----------------------------------------------------------------------------------------------------------
\section{Hamiltonian and molecular dynamics}
\label{sec:md}

The Hamiltonian is
\begin{equation}
    H[\pi,U] = \frac{1}{2}\pi^2 + S_G[U] + \sum_{k=1}^{2} \Phi_k^\dagger \bigl(DD^\dagger\bigr)^{-1} \Phi_k ,
    \label{eq:H}
\end{equation}
with $\pi^2$ as in eq.~(H.3). The pseudofermion fields are kept fixed along the trajectory, so they act as
external fields for the molecular dynamics. The leapfrog integrator (Algorithm~2 of \hmcdoc), the
accept/reject step with probability $\min[1,\exp(-\Delta H)]$, and the automatic tuning of the number of
molecular dynamics steps (Section~6 of \hmcdoc) are unchanged. Both $H[\pi',U']$ and $H[\pi,U]$ in
$\Delta H$ are evaluated with the same $\Phi_1$ and $\Phi_2$.

%----------------------------------------------------------------------------------------------------------
\section{The fermion force}
\label{sec:force}

The force is defined as in eq.~(H.24),
\begin{equation}
    F_\mu(\n) \equiv -\left.\frac{\partial S[U,\Phi_1,\Phi_2]}{\partial\omega_\mu(\n)}\right|_{\omega_\mu(\n)=0}
    = F_\mu^{\mathrm{gauge}}(\n)
    - \sum_{k=1}^{2} \Phi_k^\dagger\,
      \left.\frac{\partial (DD^\dagger)^{-1}}{\partial\omega_\mu(\n)}\right|_{\omega_\mu(\n)=0}\Phi_k ,
    \label{eq:F_def}
\end{equation}
where the gauge force is the one of eq.~(H.10),
\begin{equation}
    F_\mu^{\mathrm{gauge}}(\n) = -\beta\,\mathrm{Im}\bigl[U_\mu(\n)K_\mu^\dagger(\n)\bigr],
\end{equation}
with the staples $K_\mu(\n)$ of eq.~(H.11). The action~\eqref{eq:S} is a sum over the pseudofermion fields, so
the fermionic part of the force is the sum of the forces of the individual fields. Each term is handled exactly as
in eqs.~(H.25)--(H.29),
\begin{equation}
    \Phi_k^\dagger\,\frac{\partial (DD^\dagger)^{-1}}{\partial\omega_\mu(\n)}\,\Phi_k
    = -2\,\mathrm{Re}\left[\psi_k^\dagger\,\frac{\partial D}{\partial\omega_\mu(\n)}\,D^\dagger\psi_k\right],
    \qquad
    \psi_k = \bigl(DD^\dagger\bigr)^{-1}\Phi_k .
    \label{eq:F_k}
\end{equation}
Inserting the derivative of the Wilson operator, eqs.~(H.30) and (H.31), and defining
\begin{equation}
    \eta_k = D^\dagger\psi_k ,
    \label{eq:eta}
\end{equation}
we obtain the generalization of eq.~(H.32),
\begin{equation}
    \boxed{
    \begin{aligned}
    F_\mu(\n) = F_\mu^{\mathrm{gauge}}(\n)
    + \sum_{k=1}^{2} \mathrm{Im}\biggl(\sum_{\alpha\beta}\Bigl[
      &\psi_{k,\alpha}^*(\n)\,(1-\gamma_\mu)^{\alpha\beta}\,U_\mu(\n)\,\eta_{k,\beta}(\n+\hat\mu) \\
    - &\psi_{k,\alpha}^*(\n+\hat\mu)\,(1+\gamma_\mu)^{\alpha\beta}\,U_\mu^\dagger(\n)\,\eta_{k,\beta}(\n)
    \Bigr]\biggr).
    \end{aligned}
    }
    \label{eq:F}
\end{equation}
The gauge force enters once, while the fermionic term is the two-flavor expression evaluated once for each
pseudofermion field. Note that the field $\eta_k$ of eq.~\eqref{eq:eta} plays the role of the field called $\chi$ in
eq.~(H.32); we use a different symbol to avoid confusion with the heatbath field of eq.~\eqref{eq:heatbath}.
(On the configuration where the pseudofermions are generated, the two coincide:
$D^\dagger(DD^\dagger)^{-1}D\chi_k = \chi_k$.)

With the $\gamma$ matrices of eq.~(H.15) the spin sums can be unrolled. Writing $\psi_{k,\alpha}^*$ for the complex
conjugate of the component $\alpha$ of $\psi_k$, eqs.~(H.37) and (H.38) become
\begin{align}
    F_0(\n) - F_0^{\mathrm{gauge}}(\n) = \sum_{k=1}^{2} \mathrm{Im}\Bigl[\,
      &U_0(\n)\bigl(\psi_{k,0}^*(\n)-\psi_{k,1}^*(\n)\bigr)\bigl(\eta_{k,0}(\n+\hat0)-\eta_{k,1}(\n+\hat0)\bigr)
      \nonumber\\
    - &U_0^\dagger(\n)\bigl(\psi_{k,0}^*(\n+\hat0)+\psi_{k,1}^*(\n+\hat0)\bigr)\bigl(\eta_{k,0}(\n)+\eta_{k,1}(\n)\bigr)
      \Bigr],
    \label{eq:F0}
    \\[0.5em]
    F_1(\n) - F_1^{\mathrm{gauge}}(\n) = \sum_{k=1}^{2} \mathrm{Im}\Bigl[\,
      &U_1(\n)\bigl(\psi_{k,0}^*(\n)-i\,\psi_{k,1}^*(\n)\bigr)\bigl(\eta_{k,0}(\n+\hat1)+i\,\eta_{k,1}(\n+\hat1)\bigr)
      \nonumber\\
    - &U_1^\dagger(\n)\bigl(\psi_{k,0}^*(\n+\hat1)+i\,\psi_{k,1}^*(\n+\hat1)\bigr)\bigl(\eta_{k,0}(\n)-i\,\eta_{k,1}(\n)\bigr)
      \Bigr].
    \label{eq:F1}
\end{align}
Because of the antiperiodic boundary conditions, both terms of $F_0(\n)$ are multiplied by $-1$ when $\n$ lies on
the last time slice, since they couple $\n$ to its neighbor $\n+\hat0$ across the time boundary. In the code this is
the factor \code{rsign} returned by \code{get\_neighbors}, exactly as in the two-flavor case.

%----------------------------------------------------------------------------------------------------------
\section{Cost and alternative formulations}
\label{sec:cost}

\paragraph{Cost.}
Every evaluation of the force requires one solution of $(DD^\dagger)\psi_k = \Phi_k$ with conjugate gradient for
each pseudofermion field, followed by one application of $D^\dagger$ and one evaluation of eq.~\eqref{eq:F}. In the
leapfrog of \hmcdoc\ the force is evaluated $\NMD$ times per trajectory, and the action is evaluated twice for
$\Delta H$ (at the initial and final configurations). The number of conjugate gradient solves per trajectory is
therefore
\begin{equation}
    \Npf\,(\NMD + 2),
\end{equation}
twice as many as for two flavors at the same $\NMD$. The conjugate gradient solver itself is unchanged.

\paragraph{Step size.}
The fermion force of eq.~\eqref{eq:F} is the sum of two independent stochastic contributions, and the four-flavor
weight changes the typical gauge fields, so the integration error at a given step size differs from the two-flavor
case and $\NMD$ has to be tuned again. No manual intervention is needed: the automatic tuning of Section~6 of
\hmcdoc\ adjusts $\NMD$ during thermalization with the same target acceptance. In the test of
Section~\ref{sec:validation} ($32\times32$, $\beta=4$, $m_0=0.2$) the tuner chose the same $\NMD$ for two and four
flavors, so the cost per trajectory was dominated by the doubled number of solves.

\paragraph{Alternatives (not implemented).}
The same weight $[\det(DD^\dagger)]^2$ can be obtained with a single pseudofermion field and the action
$\Phi^\dagger(DD^\dagger)^{-2}\Phi$, generated by $\Phi = DD^\dagger\chi$. Writing $\psi = (DD^\dagger)^{-1}\Phi$
and $\xi = (DD^\dagger)^{-1}\psi$, its fermion force is
\begin{equation}
    2\,\mathrm{Re}\left[\xi^\dagger\,\frac{\partial D}{\partial\omega_\mu(\n)}\,D^\dagger\psi
    + \psi^\dagger\,\frac{\partial D}{\partial\omega_\mu(\n)}\,D^\dagger\xi\right].
\end{equation}
This also costs two conjugate gradient solves per force evaluation, but the kernel $(DD^\dagger)^{-2}$ has the
square of the condition number of $DD^\dagger$, which produces larger forces and requires a smaller step size.
Two pseudofermion fields with the kernel $(DD^\dagger)^{-1}$ are therefore preferable. Further improvements, such as
mass preconditioning~\cite{Hasenbusch2001} or the rational HMC with several pseudofermion
fields~\cite{ClarkKennedy2007} (which would also allow an odd number of flavors), would require more extensive
changes and are not part of this implementation.

%----------------------------------------------------------------------------------------------------------
\section{Implementation}
\label{sec:code}

Only the action, the force and the pseudofermion generation of the HMC depend on the number of flavors. The
Dirac operator, the conjugate gradient solver, the gauge force, the leapfrog integrator, the step-size tuner,
the halo exchange and the input/output of configurations are shared with the two-flavor code without
modification. Table~\ref{tab:code} lists the changes. Algorithm~\ref{alg:traj} summarizes one trajectory.

\begin{table}[h]
\centering
\small
\begin{tabular}{p{0.24\textwidth}p{0.70\textwidth}}
\toprule
File & Change \\
\midrule
\code{CMakeLists.txt} &
New option \code{FOUR\_FLAVORS} (default \code{OFF}), which defines the macro \code{FOUR\_FLAVORS}. Configuration
fails if it is combined with \code{CLOVER} or \code{TWISTED\_MASS}. \\
\code{include/variables.h} &
Compile-time constants \code{Nf} (4 or 2) and \code{N\_pf = Nf/2}, with an \code{\#error} if \code{FOUR\_FLAVORS} is
combined with \code{CLOVER} or \code{TWISTED\_MASS}. Fix of the copy constructor of \code{spinor\_template} (see
below). \\
\code{include/hmc.h}, \code{src/hmc.cpp} &
The pseudofermions are a \code{std::vector<spinor>} with \code{N\_pf} entries, passed to \code{Leapfrog},
\code{Force}, \code{Action} and \code{Hamiltonian}. \code{HMC\_Update} draws an independent $\chi_k$ for each field,
eq.~\eqref{eq:heatbath}. \code{Action} sums eq.~\eqref{eq:H} over the fields. \code{Force} clears the force array,
accumulates eq.~\eqref{eq:F} field by field (using the new buffer \code{Forces\_pf}), and finally adds the gauge
force. The class declares \code{HMC\_test\_access} as friend, so that the tests can call its private members. \\
\code{src/main.cpp}, \code{src/utils.cpp} &
The banner, the parameter printout and the \code{\_SimData.txt} metadata file report the number of flavors. \\
\code{run.sh}, \code{README.md} &
\code{run.sh} passes \code{-DFOUR\_FLAVORS=OFF} explicitly, like the other options, so that a cached
\code{FOUR\_FLAVORS=ON} cannot silently produce a four-flavor executable. The \code{README} describes the new build
option. \\
\code{tests/tests\_force.cpp} \newline \code{tests/reweighting/} &
Tests of Section~\ref{sec:validation}. \code{tests\_force} is a new target of the \code{BUILD\_TESTS} block of
\code{CMakeLists.txt}. \\
\bottomrule
\end{tabular}
\caption{Changes in the code for the four-flavor build.}
\label{tab:code}
\end{table}

The function \code{phi\_dag\_partialD\_phi}, which evaluates the fermionic term of eq.~\eqref{eq:F} for a single
field, overwrites its output array. For this reason \code{HMC::Force} writes the contribution of each field into
the buffer \code{Forces\_pf} and adds it to \code{Forces}. If the code is compiled with the clover term (only possible
with two flavors), the clover force is added inside the same loop, since it is also a per-field contribution.

The copy constructor of \code{spinor\_template} allocated its buffer with the member \code{size} before that member
was initialized (members are initialized in declaration order, and \code{val} is declared before \code{size}). This
did not affect the two-flavor code, which never copy-constructed a spinor, but it is triggered by the
\code{std::vector<spinor>} of pseudofermions and was fixed by allocating \code{other.size} elements.

\begin{algorithm}[h]
\caption{One HMC trajectory with $\Npf$ pseudofermion fields ($\Npf = 2$ for four flavors)}
\label{alg:traj}
\begin{tabbing}
\hspace*{1.5em}\=\hspace*{1.5em}\=\hspace*{1.5em}\=\kill
Draw $\pi_\mu(\n)$ from $\exp(-\pi^2/2)$ \quad$\triangleright$~\code{RandomPI} \\
\textbf{for} $k = 1,\dots,\Npf$ \textbf{do} \\
\> Draw $\chi_k$ from $\exp(-\chi_k^\dagger\chi_k)$ and set $\Phi_k = D[U]\,\chi_k$
 \quad$\triangleright$~\code{RandomCHI}, \code{D\_phi} \\
\textbf{end for} \\
Integrate $[\pi,U]\to[\pi',U']$ with leapfrog (Algorithm 2 of \hmcdoc), where every force evaluation is \\
\> $F \leftarrow 0$ \\
\> \textbf{for} $k = 1,\dots,\Npf$ \textbf{do} \\
\>\> $\psi_k = (DD^\dagger)^{-1}\Phi_k$, \quad $\eta_k = D^\dagger\psi_k$ \quad$\triangleright$~\code{conjugate\_gradient}, \code{D\_dagger\_phi} \\
\>\> $F \mathrel{+}=$ fermionic term of eq.~\eqref{eq:F} with $\psi_k$, $\eta_k$ \quad$\triangleright$~\code{phi\_dag\_partialD\_phi} \\
\> \textbf{end for} \\
\> $F \mathrel{+}= F^{\mathrm{gauge}}$ \quad$\triangleright$~\code{Force\_G} \\
$\Delta H = H[\pi',U'] - H[\pi,U]$ with eq.~\eqref{eq:H}, using the same $\Phi_1,\dots,\Phi_{\Npf}$ \quad$\triangleright$~\code{Hamiltonian} \\
Accept $U'$ with probability $\min[1,\exp(-\Delta H)]$
\end{tabbing}
\end{algorithm}

%----------------------------------------------------------------------------------------------------------
\section{Measurements}
\label{sec:measurements}

The program \code{mass\_NSxNT}, which computes the pion and PCAC correlators, does not depend on the number of
flavors and is unchanged. The pion correlator used there is that of a flavor non-singlet operator such as
$\bar d\gamma_5 u$. For degenerate flavors its Wick contraction contains only the connected diagram for any $\Nf$,
so the same expression applies to the four-flavor ensembles, where there are $\Nf^2-1 = 15$ degenerate pions. The
same holds for the PCAC mass. The \code{mass\_NSxNT} executable produced by the four-flavor build is identical to
the one of the two-flavor Wilson build, and either of them can be used to analyze the four-flavor configurations.

Two physical differences with respect to two flavors are worth keeping in mind when analyzing the results. First,
the critical bare mass is different. The values quoted in the \code{README} are for two flavors and must be
determined anew for $\Nf = 4$. Second, in the continuum massive $\Nf$-flavor Schwinger model the mass of the
flavor-singlet boson is $\mu^2 = \Nf g^2/\pi$, while the mass of the light non-singlet mesons scales with the
fermion mass as $m^{\Nf/(\Nf+1)}$~\cite{Coleman1976,Hetrick1995}: $m^{2/3}$ for two flavors and $m^{4/5}$ for four
flavors.

%----------------------------------------------------------------------------------------------------------
\section{Validation}
\label{sec:validation}

The implementation was tested as described below. The heatbath and force tests are in \code{tests/tests\_force.cpp}
(built with \code{-DBUILD\_TESTS=ON} and run with 4 MPI ranks). They apply to every build option. The reweighting test
is run by the script \texttt{\nolinkurl{tests/reweighting/run_reweighting_test.sh}}, which builds two- and four-flavor versions of the code on
a $4\times4$ lattice in a temporary directory.

\paragraph{Heatbath identity.}
On a random $64\times64$ gauge field, eq.~\eqref{eq:identity} holds to a relative precision of $10^{-15}$ for both
builds, with one and with $2\times2$ MPI ranks.

\paragraph{Force versus action.}
The force computed by \code{HMC::Force} was compared with a central finite difference of the action computed by
\code{HMC::Action},
\begin{equation}
    F_\mu(\n) \simeq -\frac{S\bigl[e^{i\epsilon}U_\mu(\n)\bigr] - S\bigl[e^{-i\epsilon}U_\mu(\n)\bigr]}{2\epsilon},
    \qquad \epsilon = 10^{-5},
\end{equation}
at fixed pseudofermion fields, on a random $64\times64$ gauge field with $\beta = 2$, $m_0 = 0.5$ and a conjugate
gradient tolerance of $10^{-14}$. Twelve links were tested, in both directions, including links on the last time
slice (antiperiodic boundary) and links at the boundaries between MPI subdomains. With one rank and with
$2\times2$ ranks, the largest absolute deviation is $3\times10^{-6}$, for forces of order one, which is consistent
with the solver tolerance and the rounding error of the finite difference. For comparison, the contribution of the
second pseudofermion field alone has an average magnitude of $0.85$ on the tested links, so omitting it would
be detected by more than five orders of magnitude.

\paragraph{Energy conservation.}
For an exact HMC, $\langle e^{-\Delta H}\rangle = 1$, eq.~(63) of \hmcdoc. In the runs of the next paragraph,
$\langle e^{-\Delta H}\rangle = 1.00003(16)$ and $0.99987(9)$ for four flavors, and $0.99990(10)$ and $1.00000(7)$
for two flavors.

\paragraph{Exact reweighting.}
The first three tests check the internal consistency of the force and the action, but they would not detect a wrong
weight, for instance the use of the same $\chi$ for both pseudofermion fields (Section~\ref{sec:heatbath}). The
weight itself is tested by the exact relation
\begin{equation}
    \langle O\rangle_{\Nf=4} = \frac{\bigl\langle O\,[\det(D)]^2\bigr\rangle_{\Nf=2}}{\bigl\langle [\det(D)]^2\bigr\rangle_{\Nf=2}},
    \label{eq:rw}
\end{equation}
which follows from eq.~\eqref{eq:Z}. On a $4\times4$ lattice, where $\det(D)$ can be computed exactly from the dense
$32\times32$ Wilson matrix, we generated $4\times10^5$ trajectories (four independent chains, $\NMD = 10$, trajectory
length 1) with each build, and compared the direct four-flavor averages with the right-hand side of
eq.~\eqref{eq:rw}. Table~\ref{tab:rw} shows the results for the average plaquette and for $\log|\det(D)|$, with
jackknife errors (100 blocks). The direct and reweighted averages agree within $1.1\,\sigma$, while the two- and
four-flavor values differ by more than $100\,\sigma$. In all configurations $\det(D)$ was real to machine precision,
as expected from $\gamma_5$-hermiticity. The script uses the same parameters with $2\times10^5$ trajectories per
build by default. It reports a failure if a direct and a reweighted average differ by more than $3\,\sigma$, if
$\langle e^{-\Delta H}\rangle$ differs from one by more than $3\,\sigma$, or if the two- and four-flavor averages are
not separated by at least $5\,\sigma$ (in which case the test would have no power).

\begin{table}[h]
\centering
\small
\begin{tabular}{cccccc}
\toprule
$(\beta, m_0)$ & Observable & $\Nf=2$ direct & $\Nf=4$ direct & $\Nf=2$ reweighted to $\Nf=4$ & Deviation \\
\midrule
$(2.0,\,0.0)$ & plaquette       & 0.73895(25) & 0.77673(21) & 0.77674(23) & $0.0\,\sigma$ \\
              & $\log|\det D|$  & 23.1934(12) & 23.4692(10) & 23.4703(9)  & $0.9\,\sigma$ \\
$(1.0,\,0.1)$ & plaquette       & 0.48303(33) & 0.53341(33) & 0.53397(38) & $1.1\,\sigma$ \\
              & $\log|\det D|$  & 24.1582(6)  & 24.3482(8)  & 24.3494(8)  & $1.1\,\sigma$ \\
\bottomrule
\end{tabular}
\caption{Reweighting test on a $4\times4$ lattice. The last column is the difference between the direct and the
reweighted four-flavor averages in units of the combined error.}
\label{tab:rw}
\end{table}

\paragraph{End-to-end run.}
Finally, the executables \code{SM\_32x32} of both builds were run through the standard input on $2\times2$ MPI ranks,
with $\beta = 4$, $m_0 = 0.2$, trajectory length 1, initial $\NMD = 10$, 500 thermalization trajectories with
step-size tuning and 200 measurements. Table~\ref{tab:run32} summarizes the output. The ratio of execution times,
1.8, is close to the factor 2 in the number of conjugate gradient solves at equal $\NMD$.

\begin{table}[h]
\centering
\small
\begin{tabular}{cccccc}
\toprule
$\Nf$ & Tuned $\NMD$ & Acceptance & Plaquette & Gauge action / volume & Execution time \\
\midrule
2 & 16 & 0.82 & 0.86693(48) & 0.5323(19) & 111 s \\
4 & 16 & 0.81 & 0.87117(49) & 0.5153(20) & 202 s \\
\bottomrule
\end{tabular}
\caption{Short runs on a $32\times32$ lattice with $\beta = 4$ and $m_0 = 0.2$ ($2\times2$ MPI ranks). The errors are the
jackknife errors (20 blocks) printed by the program; with 200 measurements they are only indicative.}
\label{tab:run32}
\end{table}

%----------------------------------------------------------------------------------------------------------
\begin{thebibliography}{9}

\bibitem{Hasenbusch2001}
M.~Hasenbusch, \emph{Speeding up the Hybrid Monte Carlo algorithm for dynamical fermions},
Phys.\ Lett.\ B \textbf{519} (2001) 177.

\bibitem{ClarkKennedy2007}
M.~A.~Clark and A.~D.~Kennedy, \emph{Accelerating dynamical fermion computations using the rational hybrid Monte
Carlo (RHMC) algorithm with multiple pseudofermion fields}, Phys.\ Rev.\ Lett.\ \textbf{98} (2007) 051601.

\bibitem{Coleman1976}
S.~Coleman, \emph{More about the massive Schwinger model}, Ann.\ Phys.\ \textbf{101} (1976) 239.

\bibitem{Hetrick1995}
J.~E.~Hetrick, Y.~Hosotani and S.~Iso, \emph{The massive multi-flavor Schwinger model},
Phys.\ Lett.\ B \textbf{350} (1995) 92.

\end{thebibliography}

\end{document}
